% GENERATED FILE. Edit canonical lessons, then run npm run generate:books.
% canonical-source-hash: fnv1a64:766e91bc6428dfb5
% canonical-lessons: TE-C433-krsi, TE-C433-savalu, TE-C433-vyuham, TE-C433-adhipatyam, TE-C433-kolpovu

\chapter{Effort and Strategy}
\label{ch:te-effort-and-strategy}
\hlchaptermodality{433}

\begin{chapteropening}
\textbf{By the end of this chapter:} \emph{I can talk about effort, challenges, strategy and loss.}

Complete the last lesson of chapter 433: I can talk about effort, challenges, strategy and loss.
\end{chapteropening}

\section[\texorpdfstring{\te{కృషి} (kṛṣi)}{కృషి (kṛṣi)}]{\te{కృషి} (kṛṣi) — endeavour; hard work}

\label{lesson:TE-C433-krsi}



\begin{quote}
Before the new one: say the Telugu for credibility, then the Telugu for awareness.
\end{quote}

\subsection*{You'll want to know: \te{కృషి}}
\textbf{\te{కృషి}} — \emph{kṛṣi} — \textquotedblleft{}endeavour; hard work\textquotedblright{}.

\begin{quote}
\te{ఆమె} \te{కృషి} \te{వల్లే} \te{ఆ} \te{పాఠశాల} \te{బాగుపడింది}.
\end{quote}

\begin{quote}
\emph{āme kṛṣi vallē ā pāṭhaśāla bāgupaḍindi.}
\end{quote}

\begin{quote}
That school got better because of her hard work.
\end{quote}

\te{కృషి} \te{చేయు}, to strive.

\begin{cousinweb}
\textbf{\te{కృషి}} (\emph{kṛṣi}) is borrowed from Sanskrit \emph{kṛṣi}, \textquotedblleft{}ploughing, cultivation of the soil\textquotedblright{} (Monier-Williams); Telugu uses it for patient, hard effort.
\end{cousinweb}

\begin{grammarlens}[title={What you've built}]
One more word for saying what happens if.
\end{grammarlens}

\subsection*{Your turn}
\begin{itemize}
\raggedright
  \item \emph{Say it:} \emph{kṛṣi}
  \item \emph{Say it:} \emph{kṛṣi}, once more
  \item \emph{Say it:} \emph{āme kṛṣi vallē ā pāṭhaśāla bāgupaḍindi.} — That school got better because of her hard work.
  \item \emph{Recall:} say the Telugu for credibility, then the Telugu for awareness, then say \emph{kṛṣi} again
\end{itemize}

\begin{tcolorbox}[breakable,colback=teal!4,colframe=teal!35!black,title={Before you move on}]
What does \te{కృషి} mean? (\textquotedblleft{}Endeavour; hard work\textquotedblright{}.) Say it once more.
\end{tcolorbox}
\section[\texorpdfstring{\te{సవాలు} (savālu)}{సవాలు (savālu)}]{\te{సవాలు} (savālu) — a challenge}

\label{lesson:TE-C433-savalu}



\begin{quote}
Before the new one: say the Telugu for awareness, then the Telugu for endeavour.
\end{quote}

\subsection*{You'll want to know: \te{సవాలు}}
\textbf{\te{సవాలు}} — \emph{savālu} — \textquotedblleft{}a challenge\textquotedblright{}.

\begin{quote}
\te{కొత్త} \te{ఉద్యోగం} \te{నాకు} \te{ఒక} \te{పెద్ద} \te{సవాలు}.
\end{quote}

\begin{quote}
\emph{kotta udyōgaṁ nāku oka pedda savālu.}
\end{quote}

\begin{quote}
The new job is a big challenge for me.
\end{quote}

Plural: \te{సవాళ్ళు}, \emph{savāḷḷu}.

\te{సవాలు} \te{చేయు}, to challenge someone.

\begin{cousinweb}
\textbf{\te{సవాలు}} (\emph{savālu}) is usually traced to Hindustani \emph{savāl}, \textquotedblleft{}a question, a demand\textquotedblright{}, a word of Arabic origin; in Telugu it is a challenge.
\end{cousinweb}

\begin{grammarlens}[title={What you've built}]
One more word for saying what happens if.
\end{grammarlens}

\subsection*{Your turn}
\begin{itemize}
\raggedright
  \item \emph{Say it:} \emph{savālu}
  \item \emph{Say it:} \emph{savālu}, once more
  \item \emph{Say it:} \emph{kotta udyōgaṁ nāku oka pedda savālu.} — The new job is a big challenge for me.
  \item \emph{Recall:} say the Telugu for awareness, then the Telugu for endeavour, then say \emph{savālu} again
\end{itemize}

\begin{tcolorbox}[breakable,colback=teal!4,colframe=teal!35!black,title={Before you move on}]
What does \te{సవాలు} mean? (\textquotedblleft{}A challenge\textquotedblright{}.) Say it once more.
\end{tcolorbox}
\section[\texorpdfstring{\te{వ్యూహం} (vyūhaṁ)}{వ్యూహం (vyūhaṁ)}]{\te{వ్యూహం} (vyūhaṁ) — a strategy}

\label{lesson:TE-C433-vyuham}



\begin{quote}
Before the new one: say the Telugu for endeavour, then the Telugu for a challenge.
\end{quote}

\subsection*{You'll want to know: \te{వ్యూహం}}
\textbf{\te{వ్యూహం}} — \emph{vyūhaṁ} — \textquotedblleft{}a strategy\textquotedblright{}.

\begin{quote}
\te{ఆ} \te{జట్టు} \te{కొత్త} \te{వ్యూహంతో} \te{ఆడింది}.
\end{quote}

\begin{quote}
\emph{ā jaṭṭu kotta vyūhantō āḍindi.}
\end{quote}

\begin{quote}
That team played with a new strategy.
\end{quote}

Plural: \te{వ్యూహాలు}, \emph{vyūhālu}.

\te{వ్యూహాత్మక}, strategic.

\begin{cousinweb}
\textbf{\te{వ్యూహం}} (\emph{vyūhaṁ}) is borrowed from Sanskrit \emph{vyūha}, \textquotedblleft{}arrangement\textquotedblright{}, and \textquotedblleft{}military array\textquotedblright{} (CDIAL 12222).
\end{cousinweb}

\begin{grammarlens}[title={What you've built}]
One more word for saying what happens if.
\end{grammarlens}

\subsection*{Your turn}
\begin{itemize}
\raggedright
  \item \emph{Say it:} \emph{vyūhaṁ}
  \item \emph{Say it:} \emph{vyūhaṁ}, once more
  \item \emph{Say it:} \emph{ā jaṭṭu kotta vyūhantō āḍindi.} — That team played with a new strategy.
  \item \emph{Recall:} say the Telugu for endeavour, then the Telugu for a challenge, then say \emph{vyūhaṁ} again
\end{itemize}

\begin{tcolorbox}[breakable,colback=teal!4,colframe=teal!35!black,title={Before you move on}]
What does \te{వ్యూహం} mean? (\textquotedblleft{}A strategy\textquotedblright{}.) Say it once more.
\end{tcolorbox}
\section[\texorpdfstring{\te{ఆధిపత్యం} (ādhipatyaṁ)}{ఆధిపత్యం (ādhipatyaṁ)}]{\te{ఆధిపత్యం} (ādhipatyaṁ) — dominance; supremacy}

\label{lesson:TE-C433-adhipatyam}



\begin{quote}
Before the new one: say the Telugu for a challenge, then the Telugu for a strategy.
\end{quote}

\subsection*{You'll want to know: \te{ఆధిపత్యం}}
\textbf{\te{ఆధిపత్యం}} — \emph{ādhipatyaṁ} — \textquotedblleft{}dominance; supremacy\textquotedblright{}.

\begin{quote}
\te{ఆ} \te{ఆటలో} \te{వాళ్ళ} \te{ఆధిపత్యం} \te{కొనసాగింది}.
\end{quote}

\begin{quote}
\emph{ā āṭalō vāḷḷa ādhipatyaṁ konasāgindi.}
\end{quote}

\begin{quote}
Their dominance in that game continued.
\end{quote}

\te{ఆధిపత్యం} \te{చెలాయించు}, to lord it over others.

\begin{cousinweb}
\textbf{\te{ఆధిపత్యం}} (\emph{ādhipatyaṁ}) is borrowed from Sanskrit \emph{ādhipatya}, \textquotedblleft{}supremacy, sovereignty, power\textquotedblright{} (Monier-Williams), from \emph{adhipati}, \textquotedblleft{}a ruler\textquotedblright{}.
\end{cousinweb}

\begin{grammarlens}[title={What you've built}]
One more word for saying what happens if.
\end{grammarlens}

\subsection*{Your turn}
\begin{itemize}
\raggedright
  \item \emph{Say it:} \emph{ādhipatyaṁ}
  \item \emph{Say it:} \emph{ādhipatyaṁ}, once more
  \item \emph{Say it:} \emph{ā āṭalō vāḷḷa ādhipatyaṁ konasāgindi.} — Their dominance in that game continued.
  \item \emph{Recall:} say the Telugu for a challenge, then the Telugu for a strategy, then say \emph{ādhipatyaṁ} again
\end{itemize}

\begin{tcolorbox}[breakable,colback=teal!4,colframe=teal!35!black,title={Before you move on}]
What does \te{ఆధిపత్యం} mean? (\textquotedblleft{}Dominance; supremacy\textquotedblright{}.) Say it once more.
\end{tcolorbox}
\section[\texorpdfstring{\te{కోల్పోవు} (kōlpōvu)}{కోల్పోవు (kōlpōvu)}]{\te{కోల్పోవు} (kōlpōvu) — to be deprived of; to lose}

\label{lesson:TE-C433-kolpovu}



\begin{quote}
Before the new one: say the Telugu for a strategy, then the Telugu for dominance.
\end{quote}

\subsection*{You'll want to know: \te{కోల్పోవు}}
\textbf{\te{కోల్పోవు}} — \emph{kōlpōvu} — \textquotedblleft{}to be deprived of; to lose\textquotedblright{}.

\begin{quote}
\te{వరదల్లో} \te{వాళ్ళు} \te{తమ} \te{ఇల్లు} \te{కోల్పోయారు}.
\end{quote}

\begin{quote}
\emph{varadallō vāḷḷu tama illu kōlpōyāru.}
\end{quote}

\begin{quote}
They lost their home in the floods.
\end{quote}

Past (I …): \te{కోల్పోయాను}, \emph{kōlpōyānu}.

For losing what you had: a chance, a right, a home, a person.

\begin{cousinweb}
\textbf{\te{కోల్పోవు}} (\emph{kōlpōvu}) is native Telugu, ending in \emph{pōvu}, \textquotedblleft{}to go, to be lost\textquotedblright{} (DEDR 4572); the origin of the first part is uncertain.
\end{cousinweb}

\begin{grammarlens}[title={What you've built}]
One more word for saying what happens if.
\end{grammarlens}

\subsection*{Your turn}
\begin{itemize}
\raggedright
  \item \emph{Say it:} \emph{kōlpōvu}
  \item \emph{Say it:} \emph{kōlpōvu}, once more
  \item \emph{Say it:} \emph{varadallō vāḷḷu tama illu kōlpōyāru.} — They lost their home in the floods.
  \item \emph{Recall:} say the Telugu for a strategy, then the Telugu for dominance, then say \emph{kōlpōvu} again
\end{itemize}

\begin{tcolorbox}[breakable,colback=teal!4,colframe=teal!35!black,title={Before you move on}]
What does \te{కోల్పోవు} mean? (\textquotedblleft{}To be deprived of; to lose\textquotedblright{}.) Say it once more.
\end{tcolorbox}
